\documentclass[letterpaper,11pt]{moderncv}
\moderncvstyle{banking}
\moderncvcolor{black}

\usepackage[letterpaper, margin=0.8in]{geometry}
\usepackage{xcolor}
\usepackage{fontspec}
\setmainfont{Times New Roman}
\usepackage[unicode]{hyperref}

\name{Saeyoung}{Kim}
\address{Cambridge, MA, USA}{}{}
\phone[mobile]{+1 (351) 322 5510}
\email{saeyoung\_kim@uml.edu}
\social[linkedin]{saeyoung-macx-kim}

\recipient{Hiring Team}{Anduril Industries\\Costa Mesa, CA}
\date{\today}
\opening{Dear Hiring Team,}
\closing{Sincerely,}

\begin{document}

\makelettertitle

I'm applying for the Systems Test Engineer, Space role in Costa Mesa. I've spent the last few years doing exactly this kind of work on satellite programs, and I'd like to keep doing it on systems that matter.

At Hanwha Systems, I ran AIT campaigns on military SAR satellites across EM, QM, and FM phases. I wrote and executed system-level test plans, designed the EGSE architecture, built out ETB test configurations, and chased down anomalies when hardware didn't behave during integration. I also wrote the test documentation and procedures that other engineers used, defined acceptance criteria from system-level specs, and coordinated across electrical, mechanical, and software teams to keep campaigns on track.

Before that, my PhD was in CubeSat payload development. I took a biological sensor payload from concept through cleanroom fabrication, firmware development, and system-level qualification. I built test setups, ran environmental and reliability tests, and spent a lot of time diagnosing failures at the hardware/firmware boundary. That experience is where I learned how to think about test coverage and what it actually takes to qualify flight hardware.

I should note that I'm not currently a U.S. citizen, so I'd need to discuss clearance eligibility with your team. I'm happy to be upfront about that early in the process.

Looking forward to hearing from you. Thanks for your time.

\makeletterclosing

\end{document}
